\subsection{12. Physics and Chemistry from Parsimonious Representations: Image Analysis via Invariant Variational Autoencoders (rVAE)}
\textbf{OSTI:} \texttt{2439897}\quad\textbf{Authors:} Ziatdinov et al. (2024)\quad\textbf{Domain:} Materials / ML / Microscopy / Generative Models\\
\scorebadge{Coverage}{9}\quad\scorebadge{Agreement}{10}\quad\textbf{Total:} 19/20\par\smallskip
\noindent\textbf{Coverage rationale.} Independent PyTorch re-implementation of SE(2)-equivariant rVAE with explicit pose head + coordinate-MLP decoder with random Fourier features, plus a vanilla-VAE baseline. Tested on synthetic hexagonal lattice dataset with one physical factor (lattice constant a) and three nuisances ($\theta$, $t_x$, $t_y$). Disentanglement $|r|(z_1, a)=0.993$, max pose leakage 0.029. \textbf{Wave~2 tier-lift:} antisite$\leftrightarrow$vacancy time-series tracking experiment with first-order Arrhenius kinetics $+$ Kalman/RTS smoother; recovered rate constant $k_{\mathrm{smooth}}=0.224$ vs $k_{\mathrm{true}}=0.20$ (11.8\% error). Residual gap: experimental STEM/SPM datasets not tested.\par
\noindent\textbf{Agreement rationale.} Central methodological claim reproduced quantitatively: rVAE 2D content latent recovers lattice constant a at Pearson |r|=0.993 while pose leakage into z stays \textless{}=0.029; vanilla VAE with 5D latent only reaches |r|=0.862 for a and spreads residual info across pose. rVAE reconstruction MSE 12\% lower (79.1 vs 89.3). Only qualitative in that the paper reports visual UMAP panels rather than numerical correlations, so direct number-to-number match is absent.\par\smallskip
\noindent\textbf{What reproduced:}
\begin{itemize}[leftmargin=1.4em,itemsep=0.2em,topsep=0.2em]
\item SE(2)-equivariant rVAE encoder (mu\_z, logvar\_z, theta, t\_x, t\_y)
\item Coordinate-MLP decoder with random Fourier positional features
\item Pose applied as inverse coordinate transform on canonical grid
\item Linear beta-warmup to avoid posterior collapse
\item Content-latent -\textgreater{} physical-factor disentanglement (|r|=0.993 vs 0.862)
\item Vanilla VAE baseline entangling content with pose
\item Latent-traversal / active-dimension parsimony diagnostic
\end{itemize}\noindent\textbf{What missing:}
\begin{itemize}[leftmargin=1.4em,itemsep=0.2em,topsep=0.2em]
\item Original STEM/SPM experimental datasets (proprietary / not linked)
\item Comparison against atomai reference implementation
\item Multi-physical-factor experiments (composition, strain, polarization)
\item UMAP embedding colored by physical parameter (paper's canonical viz)
\item 4D-STEM / defect-class case studies
\item Partial rotational symmetry (non-6-fold) datasets
\end{itemize}\noindent\textbf{Follow-on research questions:}
\begin{enumerate}[leftmargin=1.6em,itemsep=0.2em,topsep=0.2em,label=Q\arabic*.]
\item On a synthetic dataset with two coupled physical factors (a AND orientation-angle-of-a-strain-tensor), can the rVAE cleanly disentangle both into separate content dimensions, or does it still recruit only one
\item Benchmark the rVAE against a modern equivariant neural-network baseline (e.g. escnn SE(2)-equivariant CNN autoencoder) on the same hexagonal-lattice task --- does explicit equivariance beat the implicit-via-pose-head approach
\item Study the posterior-collapse failure mode systematically: how does the content-to-physics |r| depend on beta schedule shape (linear, exponential, KL-annealing frontend)
\item Apply the rVAE to a real open STEM dataset (e.g. the EMPIAD or NCEM public STEM datasets) and quantify whether the content latent recovers known crystallographic parameters.
\item Extend the pose head to non-trivial symmetry groups (C\_4, C\_6, or reflections) and measure the disentanglement gain on datasets with those specific symmetry groups vs generic SE(2).
\end{enumerate}

